%% main.tex — paste into Overleaf
%% Upload figures listed at the bottom of this file to Overleaf

\documentclass[sigconf,nonacm]{acmart}

\usepackage{tikz}
\usetikzlibrary{shapes, arrows.meta, positioning}
\usepackage{booktabs}
\usepackage{multirow}
\usepackage{graphicx}
\usepackage{xcolor}
\usepackage{microtype}
\usepackage{subcaption}
\usepackage{array}
\usepackage{tabularx}

\settopmatter{printacmref=false}
\renewcommand\footnotetextcopyrightpermission[1]{}
\pagestyle{plain}

% ─────────────────────────────────────────────────────────────────────────────
\title{Zero-Shot vs Few-Shot Prompting for Violence Detection using Large Language Models on Low-Resource Language (Bangla)}

\author{Sourav Saha}
\email{sourav.saha@ucf.edu}
\affiliation{%
  \institution{University of Central Florida}
  \city{Orlando, Florida, USA}
}

\begin{document}

% ─────────────────────────────────────────────────────────────────────────────
\begin{abstract}
Detecting violence-inciting content on Bangla social media is a pressing challenge with direct real-world consequences in South Asia. While prior work has benchmarked traditional NLP methods and fine-tuned transformer models on the Vio-Lens dataset, the zero-shot and few-shot capabilities of large language models (LLMs) for this task remain unexplored. This paper evaluates three LLMs — GPT-4o, LLaMA 3.1 8B, and Mistral 7B — under zero-shot and few-shot prompting conditions on Vio-Lens, a dataset of 6,046 annotated Bangla YouTube comments labeled as Direct Violence, Passive Violence, and Non-Violence. Few-shot GPT-4o achieves macro F1 of 0.720, approaching fine-tuned XLM-RoBERTa (0.729) with only 9 in-context examples and no training. Zero-shot LLaMA 3.1 8B (0.589) surpasses all n-gram baselines without training data. Mistral 7B benefits substantially from few-shot examples, with Direct Violence F1 improving from 0.137 to 0.397 — a 2.9$\times$ gain. Direct Violence remains the hardest class across all models, reflecting both class imbalance and the coded, indirect nature of inciting language in Bangla social media.
\end{abstract}

\maketitle

% ─────────────────────────────────────────────────────────────────────────────
\section{Introduction \& Problem Statement}

Social media platforms have become vehicles for inciting communal violence, particularly in South Asia. In Bangladesh and West Bengal, violence-inciting Bangla content has been directly linked to real-world incidents causing loss of life and property. Automated detection of such content is therefore a critical public safety problem.

Despite substantial progress in hate speech detection for high-resource languages~\cite{nobata2016,bdshs}, Bangla remains severely underrepresented in NLP research. The Vio-Lens dataset~\cite{violens} addresses this gap with 6,046 annotated Bangla YouTube comments, with prior work establishing strong n-gram SVM and fine-tuned transformer baselines including BanglaBERT (macro F1 = 0.788). However, the capabilities of instruction-following LLMs for Bangla violence detection without any task-specific training remain an open question.

This is especially relevant given the practical constraints of low-resource settings: annotating training data in Bangla is expensive, models require retraining as language evolves, and fine-tuning large transformers demands significant compute. LLMs offer a compelling alternative if their zero-shot or few-shot performance is competitive.

This paper investigates three questions: (1) Can zero-shot LLM prompting match or exceed n-gram baselines without any training data? (2) Does few-shot prompting with labeled in-context examples significantly improve performance? (3) How do GPT-4o, LLaMA 3.1 8B, and Mistral 7B compare on this task?

Our key contributions are: (i) the first systematic LLM evaluation on Bangla violence detection across three models and two prompting conditions; (ii) a structured few-shot prompting strategy using a four-step decision framework grounded in Galtung's theory of violence~\cite{galtung}; and (iii) a detailed per-class error analysis revealing model-specific failure patterns on Direct Violence, the most critical and difficult class.

% ─────────────────────────────────────────────────────────────────────────────
\section{Related Work}

\paragraph{Violence and Hate Speech Detection.}
Online hate speech detection has been extensively studied for English and other high-resource languages~\cite{nobata2016,silva2016}. Bangla-specific efforts include BD-SHS~\cite{bdshs} for abusive language detection and SentNob~\cite{sentnob} for sentiment analysis on noisy Bangla text. Vio-Lens~\cite{violens} is the first dataset targeting communal violence incitement in Bangla specifically, providing an annotation framework grounded in Galtung's theory of violence~\cite{galtung}. Saha et al.~\cite{violens} establish strong baselines using n-gram SVMs and fine-tuned transformer models, with BanglaBERT achieving macro F1 of 0.788 on the validation set.

\paragraph{LLMs for Classification.}
GPT-4 and instruction-tuned LLMs have demonstrated strong zero-shot and few-shot performance across diverse NLP tasks~\cite{brown2020,openai2023gpt4}. However, their effectiveness on low-resource and sensitive tasks — particularly violence detection in non-English languages — remains largely unexplored. Our work directly addresses this gap for Bangla.

\paragraph{Prompting Strategies.}
Zero-shot prompting relies solely on task descriptions in the system prompt, while few-shot prompting augments with labeled in-context examples~\cite{brown2020}. Few-shot prompting has been shown to substantially improve performance on classification tasks, particularly for minority classes. We compare both strategies under identical evaluation conditions across three models.

% ─────────────────────────────────────────────────────────────────────────────
\section{Methodology}

\subsection{Dataset}

We use Vio-Lens~\cite{violens}, comprising 6,046 annotated Bangla YouTube comments across three labels. Table~\ref{tab:dataset} summarizes the dataset statistics. The class distribution reflects real-world imbalance: Non-Violence (NV, 51.4\%), Passive Violence (PV, 34.1\%), and Direct Violence (DV, 14.4\%), making macro-averaged F1 the appropriate primary metric. All LLM experiments are evaluated on the 1,330-sample validation split for direct comparison with published baselines~\cite{violens}.

\begin{table}[t]
\caption{Vio-Lens dataset split statistics.}
\label{tab:dataset}
\setlength{\tabcolsep}{5pt}
\begin{tabular}{lrrrr}
\toprule
\textbf{Split} & \textbf{NV} & \textbf{PV} & \textbf{DV} & \textbf{Total} \\
\midrule
Train      & 1,388 & 921  & 391 & 2,700 \\
Validation & 717   & 417  & 196 & 1,330 \\
Test       & 1,042 & 686  & 288 & 2,016 \\
\midrule
\textbf{Total} & 3,147 & 2,024 & 875 & \textbf{6,046} \\
\bottomrule
\end{tabular}
\end{table}

Table~\ref{tab:examples} shows representative examples from each class. Direct Violence posts often rely on religious or political framing rather than explicit vocabulary, which contributes to the classification difficulty observed across all models.

\begin{table}[t]
\caption{Representative examples from Vio-Lens (translated to English).}
\label{tab:examples}
\small
\begin{tabularx}{\columnwidth}{@{}p{0.55\columnwidth}X@{}}
\toprule
\textbf{Bangla Text (Excerpt)} & \textbf{Label} \\
\midrule
``I want the truth to come out. Let proper punishment be given.'' & NV \\
\midrule
``Not everyone gets the fortune of killing atheists and walking to the gallows smiling.'' & PV \\
\midrule
``Rise for Jihad, Rise for Jihad'' / ``India should deport all these terrorists to Pakistan or Afghanistan!'' & DV \\
\bottomrule
\end{tabularx}
\end{table}

\subsection{Models and Infrastructure}

We evaluate three LLMs:
\begin{itemize}[leftmargin=*,topsep=2pt,itemsep=1pt]
  \item \textbf{GPT-4o}~\cite{openai2023gpt4}: evaluated via the OpenAI API under both zero-shot and few-shot conditions.
  \item \textbf{LLaMA 3.1 8B Instruct}: evaluated zero-shot via the Groq inference API (\texttt{llama-3.1-8b-instant}), equivalent to Meta's \texttt{Meta-Llama-3.1-8B-Instruct}.
  \item \textbf{Mistral 7B Instruct v0.2}~\cite{mistral7b}: run locally on 4$\times$ NVIDIA RTX A4500 (84\,GB total VRAM) under both conditions.
\end{itemize}

All generations use greedy decoding (temperature = 0) for reproducibility, with \texttt{max\_tokens = 5} to constrain output to a single classification digit.

\subsection{Prompting Strategies}

Figure~\ref{fig:pipeline} illustrates the inference pipeline. Table~\ref{tab:prompts} compares the two prompting strategies.

\paragraph{Zero-Shot.}
The system prompt encodes precise category definitions for labels 0, 1, and 2 along with strict output rules (single digit only). The user turn contains the raw Bangla text. No examples are provided.

\paragraph{Few-Shot.}
The system prompt is augmented with a four-step decision framework grounded in Galtung's theory of violence~\cite{galtung}. This framework encodes the annotation guidelines from~\cite{violens} as a structured decision tree: it first asks whether the post justifies violence (Step 1), whether violence is direct or passive (Step 2--3), and defaults to Non-Violence otherwise (Step 4). We sample 9 labeled examples from the training set (3 per class) as interleaved user/assistant turns, forming in-context demonstrations prior to the test query. For Mistral 7B, which does not support a system role, the system prompt is prepended to the first user turn.

\begin{table}[t]
\caption{Prompt structure comparison.}
\label{tab:prompts}
\small
\setlength{\tabcolsep}{4pt}
\begin{tabularx}{\columnwidth}{@{}lXX@{}}
\toprule
 & \textbf{Zero-Shot} & \textbf{Few-Shot} \\
\midrule
\textbf{System} & Category definitions (0/1/2) + output rules & 4-step decision framework + definitions + output rules \\
\midrule
\textbf{Context} & — & 9 user/asst turns (3 per class) from train set \\
\midrule
\textbf{User} & ``Classify this Bangla text: [T]'' & ``Classify this Bangla text: [T]'' \\
\midrule
\textbf{Output} & Single digit: \texttt{0}, \texttt{1}, or \texttt{2} & Single digit: \texttt{0}, \texttt{1}, or \texttt{2} \\
\bottomrule
\end{tabularx}
\end{table}

\begin{figure}[t]
\centering
\begin{tikzpicture}[
  node distance=0.42cm,
  box/.style={rectangle, rounded corners=3pt, draw=black!65, fill=blue!7,
              text width=3.6cm, align=center, minimum height=0.62cm, font=\small},
  side/.style={rectangle, rounded corners=2pt, draw=orange!70, fill=orange!8,
               text width=2.0cm, align=center, minimum height=0.55cm, font=\scriptsize, dashed},
  arrow/.style={-Stealth, thick, draw=black!55}
]
  \node[box] (input)  {Vio-Lens Validation\\(Bangla text)};
  \node[box, below=of input]  (build)  {Prompt Builder\\{\footnotesize sys prompt + query}};
  \node[box, below=of build]  (llm)    {LLM\\{\footnotesize GPT-4o / LLaMA 3.1 / Mistral}};
  \node[box, below=of llm]    (parse)  {Response Parser\\{\footnotesize extract digit 0/1/2}};
  \node[box, below=of parse]  (out)    {Predicted Label\\{\footnotesize 0=NV · 1=PV · 2=DV}};
  \node[side, right=0.55cm of build] (pool)
       {Few-Shot Pool\\9 examples\\(3 per class)};
  \draw[arrow] (input) -- (build);
  \draw[arrow] (build) -- (llm);
  \draw[arrow] (llm)   -- (parse);
  \draw[arrow] (parse) -- (out);
  \draw[arrow, dashed, draw=orange!60] (pool) -- (build);
  \node[font=\tiny, text=gray, below=0.05cm of pool] {few-shot only};
\end{tikzpicture}
\caption{Inference pipeline. Zero-shot omits the few-shot pool; few-shot injects 9 labeled examples as interleaved user/assistant turns before the test query.}
\label{fig:pipeline}
\end{figure}

% ─────────────────────────────────────────────────────────────────────────────
\section{Evaluation \& Results}

\subsection{Main Results}

Table~\ref{tab:results} presents per-class and macro F1 on the Vio-Lens validation set. Published baseline results are taken directly from~\cite{violens}.

\begin{table*}[t]
\caption{Per-class precision, recall, and F1 on the Vio-Lens validation set. Baseline results from Saha et al.~\cite{violens}. \textbf{Bold} = best in group.}
\label{tab:results}
\setlength{\tabcolsep}{6pt}
\begin{tabular}{llcccccccc}
\toprule
 & \textbf{Model} & \multicolumn{2}{c}{\textbf{Direct Violence}} & \multicolumn{2}{c}{\textbf{Passive Violence}} & \multicolumn{2}{c}{\textbf{Non-Violence}} & \textbf{Macro F1} \\
\cmidrule(lr){3-4}\cmidrule(lr){5-6}\cmidrule(lr){7-8}
 & & P & F1 & P & F1 & P & F1 & \\
\midrule
\multirow{3}{*}{\rotatebox[origin=c]{90}{\small N-gram}}
 & Unigram SVM   & —    & 0.444 & —    & 0.629 & —    & 0.776 & 0.616 \\
 & Char-2-gram   & —    & 0.490 & —    & 0.630 & —    & 0.788 & 0.636 \\
 & Char-3-gram   & —    & 0.505 & —    & 0.662 & —    & 0.803 & 0.657 \\
\midrule
\multirow{4}{*}{\rotatebox[origin=c]{90}{\small Fine-tuned}}
 & DistilBERT    & —    & 0.548 & —    & 0.637 & —    & 0.771 & 0.652 \\
 & MBERT         & —    & 0.601 & —    & 0.643 & —    & 0.802 & 0.682 \\
 & XLM-RoBERTa  & —    & 0.670 & —    & 0.704 & —    & 0.729 & 0.729 \\
 & BanglaBERT    & —    & \textbf{0.754} & —    & \textbf{0.764} & —    & \textbf{0.846} & \textbf{0.788} \\
\midrule
\multirow{5}{*}{\rotatebox[origin=c]{90}{\small LLM (ours)}}
 & Zero-Shot GPT-4o       & 0.52 & 0.512 & 0.61 & 0.689 & 0.84 & 0.767 & 0.656 \\
 & Few-Shot GPT-4o        & 0.90 & \textbf{0.625} & 0.63 & 0.711 & 0.84 & 0.825 & \textbf{0.720} \\
 & Zero-Shot LLaMA 3.1 8B & 0.40 & 0.442 & 0.61 & 0.602 & 0.75 & 0.724 & 0.589 \\
 & Zero-Shot Mistral 7B   & 0.65 & 0.137 & 0.46 & 0.446 & 0.64 & 0.711 & 0.431 \\
 & Few-Shot Mistral 7B    & 0.41 & 0.397 & 0.46 & 0.532 & 0.76 & 0.664 & 0.531 \\
\bottomrule
\end{tabular}
\end{table*}

\paragraph{Zero-Shot Performance.}
Zero-shot GPT-4o achieves macro F1 of 0.656, surpassing all n-gram SVM baselines and falling only 0.026 points below fine-tuned MBERT — without any training data. Zero-shot LLaMA 3.1 8B (0.589) outperforms Unigram and Char-2-gram SVM baselines. Zero-shot Mistral 7B (0.431) lags significantly, particularly on Direct Violence (F1 = 0.137), where it produces near-degenerate predictions by rarely assigning the minority label.

\paragraph{Few-Shot Performance.}
Few-shot GPT-4o reaches macro F1 of 0.720, a gain of +0.064 over zero-shot, approaching XLM-RoBERTa (0.729). The gain is most pronounced on Direct Violence (+0.113 F1), the hardest class. Few-shot Mistral 7B (0.531) improves substantially over its zero-shot baseline (+0.100), with Direct Violence F1 jumping from 0.137 to 0.397 — the largest within-model gain observed. Three labeled examples per class are sufficient to calibrate Mistral's decision boundary on the minority class.

\paragraph{Cross-Model Comparison.}
GPT-4o consistently outperforms both open-weight models across all conditions. Zero-shot LLaMA 3.1 8B (0.589) outperforms Mistral 7B under both zero-shot (0.431) and few-shot (0.531), suggesting that LLaMA's stronger multilingual instruction-following capability matters more for Bangla than Mistral's access to few-shot demonstrations at this scale.

\begin{figure*}[t]
\centering
\includegraphics[width=0.92\textwidth]{macro_f1_comparison.png}
\caption{Macro F1 comparison across all models. Dashed line separates supervised baselines (left) from LLM prompting (right). Few-shot GPT-4o approaches XLM-RoBERTa with no training.}
\label{fig:macro}
\end{figure*}

\subsection{Per-Class Analysis}

Figure~\ref{fig:perclass} shows per-class F1 across all five LLM conditions.

Non-Violence F1 is consistently the highest across all models (0.664--0.825), reflecting the majority class advantage. Passive Violence is moderate across all models, with a tendency toward over-prediction due to lexical overlap with Direct Violence posts. Direct Violence is the hardest class in every condition: even BanglaBERT achieves only F1 = 0.754, and few-shot GPT-4o reaches 0.625. The challenge stems from class imbalance (14.4\% of samples), and the fact that many violent posts use indirect, metaphorical, or coded language.

\begin{figure*}[t]
\centering
\includegraphics[width=0.92\textwidth]{per_class_f1.png}
\caption{Per-class F1 by model and prompting condition. Direct Violence (red) is consistently the hardest class. Few-shot examples provide the largest gains for DV across both models.}
\label{fig:perclass}
\end{figure*}

\subsection{Confusion Matrices}

Figure~\ref{fig:cm} shows confusion matrices for GPT-4o and Mistral 7B under both conditions. Several patterns emerge. Zero-shot Mistral 7B systematically misclassifies Direct Violence as Non-Violence (only 16 correct out of 196 DV samples), consistent with its DV F1 of 0.137. Few-shot exposure corrects this substantially. GPT-4o few-shot achieves high Direct Violence precision (0.90) — when it predicts DV, it is almost always correct — but recall remains low (0.48), indicating a conservative prediction tendency for the minority class.

\begin{figure*}[t]
\centering
\includegraphics[width=0.95\textwidth]{confusion_matrices_4panel.png}
\caption{Confusion matrices for GPT-4o (top) and Mistral 7B (bottom) under zero-shot (left) and few-shot (right) conditions. Note the dramatic improvement in Direct Violence identification with few-shot prompting, especially for Mistral.}
\label{fig:cm}
\end{figure*}

\subsection{Zero-Shot vs.\ Few-Shot Gains}

Figure~\ref{fig:delta} shows the per-class F1 gain from adding few-shot examples for GPT-4o and Mistral. Both models benefit most on Direct Violence. GPT-4o gains +0.113 on DV and +0.055 on Non-Violence, while Mistral gains +0.260 on DV and +0.086 on Passive Violence. Notably, Mistral's Non-Violence F1 declines slightly (−0.047) with few-shot, suggesting that examples reshape the decision boundary toward Direct Violence at the cost of some Non-Violence recall.

\begin{figure*}[t]
\centering
\includegraphics[width=0.88\textwidth]{zeroshot_vs_fewshot_delta.png}
\caption{Zero-shot vs.\ few-shot F1 comparison for GPT-4o (left) and Mistral 7B (right). Delta values shown above each bar pair. Both models gain most on Direct Violence, the hardest class.}
\label{fig:delta}
\end{figure*}

% ─────────────────────────────────────────────────────────────────────────────
\section{Discussion \& Challenges}

\paragraph{LLM Capability Gap.}
The gap between GPT-4o and the open-weight models is substantial and consistent. GPT-4o's advantage likely stems from larger-scale multilingual pretraining and superior instruction following in non-English languages. The fact that zero-shot LLaMA 3.1 8B (0.589) outperforms even few-shot Mistral 7B (0.531) suggests that base capability dominates over prompting strategy at smaller scales for Bangla — a finding with practical implications for resource-constrained deployments.

\paragraph{The Direct Violence Problem.}
Direct Violence is the most challenging class across all conditions. Zero-shot Mistral's F1 of 0.137 indicates near-random performance, driven by systematic avoidance of the minority label. Yet three in-context DV examples triple this score to 0.397, demonstrating that labeled examples are an efficient calibration mechanism for minority class identification even when base zero-shot performance is very low. GPT-4o few-shot achieves the highest DV precision (0.90) but suffers from low recall (0.48), suggesting the model learns to be very confident when it does predict DV but still misses many cases.

\paragraph{Infrastructure Challenges.}
Running LLMs at research scale poses practical challenges. Groq's free tier (500k tokens/day) was sufficient for zero-shot LLaMA (1,330 samples $\times$ $\sim$200 tokens), but insufficient for few-shot (1,330 $\times$ $\sim$1,500 tokens). Mistral 7B was run locally on a 4$\times$ A4500 GPU cluster, with per-sample inference taking $\sim$3--4 seconds under few-shot. These constraints are directly relevant for reproducibility and accessibility in low-resource research settings.

\paragraph{Prompt Sensitivity.}
Minor variations in category definitions significantly affect Mistral's output distribution. The structured four-step decision framework substantially improves Mistral's DV precision by providing explicit reasoning guidance that the model cannot derive from definitions alone. GPT-4o is more robust to prompt variations, consistently applying criteria without explicit decision structure.

\paragraph{Limitations.}
LLaMA 3.1 8B few-shot results were unavailable due to daily API token limits; only zero-shot is reported. All experiments use the validation split; generalization to the held-out test set is unknown. The few-shot pool is fixed at 9 examples and sensitivity to example selection was not evaluated.

% ─────────────────────────────────────────────────────────────────────────────
\section{Concluding Remarks}

This paper presents the first systematic evaluation of LLM prompting for Bangla violence detection, covering GPT-4o, LLaMA 3.1 8B, and Mistral 7B under zero-shot and few-shot conditions on the Vio-Lens benchmark. Few-shot GPT-4o achieves macro F1 of 0.720, approaching fine-tuned XLM-RoBERTa (0.729) with only 9 labeled examples. Zero-shot LLaMA 3.1 8B outperforms all n-gram baselines without any training. Few-shot examples consistently improve Direct Violence detection — the hardest and most safety-critical class — across all models, with Mistral showing a 2.9$\times$ DV F1 gain.

Future work will evaluate LLaMA 3.1 8B under few-shot conditions, investigate robustness to semantic-preserving perturbations (paraphrasing, word masking, character noise), and explore retrieval-augmented example selection to improve minority class detection beyond fixed pools.

% ─────────────────────────────────────────────────────────────────────────────
\bibliographystyle{ACM-Reference-Format}
\bibliography{refs}

\end{document}
